\documentclass[10pt]{article}
\usepackage[margin=1.7cm]{geometry}
\usepackage{amsmath,amssymb,amsthm}
\pagestyle{empty}

\newtheorem{problem}{Problem}
\newtheorem{lemma}{Lemma}
\theoremstyle{definition}
\newtheorem{answer}{Answer}

\title{\vspace{-1.4em}\textbf{A Digit-Concatenation Puzzle}\\[0.2em]
\large Complete Solution}
\author{}
\date{}

\begin{document}
\maketitle
\vspace{-4.2em}

\begin{problem}
Let $A$ be a two-digit number and $B$ a three-digit number such that the product
$C=A\cdot B$ is a five-digit number and the ten digits of $A,B,C$ taken together are
exactly $0,1,\dots,9$ (each used once). Suppose furthermore that
\[
A<R(A),\qquad B>R(B),\qquad B+R(B)=11\bigl(A+R(A)\bigr),
\]
where $R(x)$ denotes the integer obtained by reversing the digits of $x$.
Determine the ten-digit number $N$ formed by concatenating $A$, $B$, and $C$.
\end{problem}

\section*{Notation and First Reductions}
Write
\[
A=10a+b,\qquad B=100c+10d+e,
\]
where $a,b,c,d,e$ are pairwise distinct digits and $a,c\neq 0$. Reversing gives
\[
R(A)=10b+a,\qquad R(B)=100e+10d+c.
\]
The inequalities reduce to
\[
A<R(A)\iff 10a+b<10b+a\iff a<b,
\qquad
B>R(B)\iff 100c+10d+e>100e+10d+c\iff c>e.
\]
Put $s=a+b$. Then
\[
A+R(A)=11(a+b)=11s,
\]
and the given reversal relation $B+R(B)=11(A+R(A))$ becomes
\[
B+R(B)=101(c+e)+20d=121s. \tag{1}
\]

\subsection*{Two congruences from (1)}
Reducing (1) modulo $10$ gives
\[
c+e\equiv s\pmod{10}. \tag{2}
\]
Reducing (1) modulo $11$, using $101\equiv 2$, $20\equiv 9$, $121\equiv 0\pmod{11}$,
\[
2(c+e)+9d\equiv 0\pmod{11}
\iff c+e\equiv d\pmod{11}. \tag{3}
\]

\begin{lemma}\label{lem:core}
$a+b=d=c+e$.
\end{lemma}
\begin{proof}
Since $0\le c+e\le 17$ and $0\le d\le 9$, congruence (3) leaves only
\[
c+e=d \quad\text{or}\quad c+e=d+11.
\]
If $c+e=d+11$, substituting into (1) gives
\[
101(d+11)+20d=121s
\iff 121d+1111=121s
\iff 1111=121(s-d),
\]
impossible, because $1111\equiv 22\pmod{121}$ is not divisible by $121$.
Hence $c+e=d$, and (1) becomes $121d=121s$, so $s=d$.
\end{proof}

Consequently all six digits $a,b,c,d,e$ and~$0$ are forced to be distinct
where they occur: indeed $e\neq 0$ (otherwise $c+e=d$ forces $c=d$, contradicting
distinctness), $b\neq 0$ since $a<b$, and $a,c\neq 0$ by hypothesis. Therefore the
digit $0$ must appear in $C$.

\section*{Modulo 9: Narrowing the Middle Digit}
The digit sums are
\[
A:\ a+b=d,\qquad B:\ c+d+e=(c+e)+d=2d,\qquad C:\ 45-3d,
\]
because all ten digits occur exactly once. Since $C=A\cdot B$ and every number is
congruent to the sum of its digits modulo $9$,
\[
45-3d\equiv d\cdot 2d=2d^{2}\pmod 9
\iff d(2d+3)\equiv 0\pmod 9.
\]
As $d=c+e$ with $c>e\ge 1$ and $d\le 9$, we have $3\le d\le 9$. Checking
$d=3,\dots,9$:
\[
d=3:\ 3\cdot 9=27\equiv 0,\qquad
d=6:\ 6\cdot 15=90\equiv 0,\qquad
d=9:\ 9\cdot 21=189\equiv 0,
\]
while $d=4,5,7,8$ give residues $8,2,2,8$. Hence
\[
d\in\{3,6,9\}.
\]

\section*{Case Analysis}
\subsection*{Case $d=3$.}
$c+e=3$ with $c>e\ge 1$ forces $(c,e)=(2,1)$, so $B=231$. Then $a+b=3$ with
$a<b$ forces $A=12$, but the digits $1,2$ already occur in $B$. No solution.

\subsection*{Case $d=6$.}
$c+e=6$ with $c>e\ge 1$: $(c,e)=(5,1)$ or $(4,2)$.
\begin{itemize}
\item $B=561$: $a+b=6$ with $a<b$ and digits distinct from $\{5,6,1\}$ leaves only
$A=24$ (the pairs $(0,6)$ and $(3,3)$ are excluded by $a\neq 0$ and distinctness).
But $24\cdot 561=13464$ has repeated digits. Rejected.
\item $B=462$: $a+b=6$ leaves $A=15$, but $15\cdot 462=6930$ is not a five-digit
number. Rejected.
\end{itemize}

\subsection*{Case $d=9$.}
$c+e=9$ with $c>e\ge 1$: $(c,e)\in\{(8,1),(7,2),(6,3),(5,4)\}$, so
\[
B\in\{891,792,693,594\},
\]
and $a+b=9$ with $a<b$ gives
\[
A\in\{18,27,36,45\}.
\]
Discard candidates by two tests: (i) any digit of $A$ already occurring in $B$;
(ii) the units digit of $C$, which is $be \bmod 10$, colliding with a used digit.

\begin{center}
\footnotesize
\begin{tabular}{p{1.0cm}|p{2.5cm}|p{10.8cm}}
$B$ & $A$ after test (i) & outcome \\ \hline
$891$ & $27,36,45$ & units digit of $C$ equals $b\in\{7,6,5\}$, already used --- all fail \\ \hline
$792$ & $18,36,45$ & $18\cdot 792=14256$ repeats $1,2$; $36$: units $6\cdot 2\equiv 2$ used; $45\cdot 792=35640$ repeats $4,5$ --- all fail \\ \hline
$693$ & $18,27,45$ & $18\cdot 693=12474$ has repeated digits; $27\cdot 693=18711$ repeats; $45$: units $5\cdot 3\equiv 5$ used --- all fail \\ \hline
$594$ & $18,27,36$ & $18\cdot 594=10692$ repeats $1,9$; $36$: units $6\cdot 4\equiv 4$ used; $\mathbf{27\cdot 594=16038}$ works
\end{tabular}
\end{center}

\section*{Verification and Answer}
\[
27\cdot 594=16038,\qquad 27<72,\qquad 594>495,
\]
\[
594+495=1089=11\cdot 99=11(27+72).
\]
All conditions hold, and
\[
N=27\,|\,594\,|\,16038=2759416038
\]
uses every digit $0$--$9$ exactly once.

\begin{answer}
\[
\boxed{\;A=27,\quad B=594,\quad C=16038,\quad N=2759416038.\;}
\]
\end{answer}

\end{document}
